%!TEX root = ../../../note

\subsection{Numbers and Logic}

\subsubsection{What Is Calculus?}
Differential and integral calculus are two branches of calculus.

\begin{description}
    \item[Differential calculus:] The study of slopes and tangent lines.
    \item[Integral calculus:] The study of areas and accumulation.
\end{description}

\subsubsection{From $\N$ to $\Q$}
The number systems discussed in this lecture include
\begin{align*}
    \N &= \set{1,2,3,\ldots}, \\
    \Z &= \set{\ldots,-1,0,1,\ldots}, \\
    \Q &= \set{\frac{m}{n} : m,n\in\Z,\ n\ne 0}.
\end{align*}

The natural numbers are closed under addition. The integers are closed under addition, subtraction, and multiplication, but they are not closed under division; for example, $1/2\notin\Z$. Therefore, we introduce the rational numbers to include quotients of integers with nonzero denominators.

\begin{theorem}[Irrationality of $\sqrt{2}$]
    The number $\sqrt{2}$ is irrational; equivalently, $\sqrt{2}\notin\Q$.
    \insymbols{\neg(\exists m\in\Z)(\exists n\in\N)\,\Bigl[\sqrt2=\frac mn\Bigr]}
\end{theorem}

\begin{idea}
    Write $\sqrt2=m/n$ in lowest terms. Then $m^2=2n^2$ forces $m$ even, and feeding that back forces $n$ even --- so $2$ divides both, contradicting lowest terms.
\end{idea}

\begin{proof}
    Suppose, for contradiction, that $\sqrt{2}\in\Q$. Then we can write
    \begin{equation*}
        \sqrt{2}=\frac{m}{n},
        \qquad m,n\in\Z,
        \qquad n>0,
        \qquad \gcd(m,n)=1.
    \end{equation*}
    Squaring both sides gives
    \begin{equation*}
        2=\frac{m^2}{n^2}
        \qquad\Longrightarrow\qquad
        m^2=2n^2.
    \end{equation*}

    Thus $m^2$ is even. An odd integer has an odd square, so $m$ must be even. Write $m=2\ell$ for some $\ell\in\Z$. Substituting this into the equation above gives
    \begin{align*}
        2n^2 &= 4\ell^2, \\
        n^2 &= 2\ell^2.
    \end{align*}
    Therefore $n^2$ is even, and the same parity argument shows that $n$ is even.

    This means that both $m$ and $n$ are even, contradicting $\gcd(m,n)=1$. Therefore, $\sqrt{2}\notin\Q$.
\end{proof}

\[\boxed{\text{Assuming lowest terms turns the equation into a parity contradiction.}}\]

The irrational numbers are the real numbers that are not rational:
\begin{equation*}
    \R\setminus\Q
    =\set{x\in\R : x\notin\Q}.
\end{equation*}

\begin{example}[Decimal representation]
Every nonnegative real number has a decimal representation of the form
\begin{equation*}
    x=a+\sum_{i=1}^{\infty}\frac{x_i}{10^i},
    \qquad a\in\N\cup\set{0},
    \qquad x_i\in\set{0,1,2,\ldots,9}.
\end{equation*}
Negative real numbers are the negatives of such expansions. Some numbers have two decimal representations, for example $1.999\cdots=2.000\cdots$.
\end{example}

\begin{example}[Recurring decimal]
    Is the statement $1.999\cdots = 2$ correct?

    Yes. Let $x=1.999\cdots$. Then
    \begin{align*}
        x &= 1.999\cdots, \\
        10x &= 19.999\cdots, \\
        9x &= 18, \\
        x &= 2.
    \end{align*}
\end{example}

\subsubsection{Statements and Quantifiers}
\begin{definition}[Statement]
    A statement is a sentence that has a definite truth value: it is either true or false.
\end{definition}

The \textbf{universal quantifier} $\forall$ means ``for all.'' For example,
\begin{example}
    The statement
    \begin{equation*}
        (\forall x \in \R)\,[x^2-x\geq 0]
    \end{equation*}
    is false; for instance, it fails when $x=\frac12$.
\end{example}

The \textbf{existential quantifier} $\exists$ means ``there exists.'' For example,
\begin{example}
    The statement
    \begin{equation*}
        (\exists x \in \R)\,[x^2-x\geq 0]
    \end{equation*}
    is true; for instance, $x=2$ satisfies the inequality.
\end{example}

\begin{notation}
    Each quantifier is written in parentheses, and the statement it governs in square brackets, read left to right:
    \[
        (\forall x\in\R)(\exists y\in\R)\,[x^3-y^3=1]
    \]
    reads ``for every $x\in\R$ there exists $y\in\R$ such that $x^3-y^3=1$.'' Negation moves inward and swaps each quantifier:
    \[
        \neg(\forall x)(\exists y)\,[P(x,y)] \iff (\exists x)(\forall y)\,[\neg P(x,y)].
    \]
    From here on, every theorem-like statement ends with its form \emph{In symbols}: quantifiers and connectives only, using notation defined before that statement.
\end{notation}

\begin{proposition}
    For every $x\in\R$, there exists $y\in\R$ such that
    \begin{equation*}
        x^3-y^3=1.
    \end{equation*}
    \insymbols{(\forall x\in\R)(\exists y\in\R)\,[x^3-y^3=1]}
\end{proposition}

\begin{proof}
    Let $x\in\R$ and take $y=\sqrt[3]{x^3-1}$. Then
    \begin{equation*}
        x^3-y^3=x^3-(x^3-1)=1.
    \end{equation*}
\end{proof}

Here the choice of $y$ is allowed to depend on $x$: for each $x$ we produce a (possibly different) $y$ that makes the equation hold. This is exactly what the quantifier order $(\forall x)(\exists y)$ permits.

\begin{proposition}
    The statement
    \begin{equation*}
        (\exists y\in\R)(\forall x\in\R)\,[x^3-y^3=1]
    \end{equation*}
    is false.
    \insymbols{\neg(\exists y\in\R)(\forall x\in\R)\,[x^3-y^3=1]}

    Its negation is
    \begin{equation*}
        (\forall y\in\R)(\exists x\in\R)\,[x^3-y^3\neq 1].
    \end{equation*}
    To prove the negation, let $y\in\R$ and take $x=y$. Then
    $x^3-y^3=0\neq 1$.
\end{proposition}

\begin{remark}
    In this second statement, the quantifier order is reversed: $(\exists y)(\forall x)$ demands a \emph{single} $y$ that works simultaneously for every $x\in\R$. This is a much stronger requirement than $(\forall x)(\exists y)$, which only asks that a (possibly different) $y$ exist for each individual $x$. The first proposition shows such $y$'s exist, but they genuinely depend on $x$ (namely $y=\sqrt[3]{x^3-1}$), so no single choice of $y$ can work for all $x$ at once. This is why swapping the order of $\forall$ and $\exists$ can change a true statement into a false one, and it is a caution to always read quantified statements in the order they are written.
\end{remark}
